\begin{frame}{스케일링은 답을 바꾸는 게 아니다}
  \begin{itemize}
    \item 정규방정식은 두 경우 모두 \textbf{한 번의 행렬 계산으로
      정확히} $(b = 20, \text{slope} = 1.5)$ 를 낸다.
    \item 표준화 케이스는 $b = 155$, $\text{slope}_s = 42.43$ 이
      나와도, 스케일 변환으로 역변환하면
      $\text{slope} = 42.43 / \sigma = 1.5$ 가 된다.
  \end{itemize}
  \vskip0.8em
  \begin{block}{이 절에서 꼭 잡아야 할 핵심}
    \kb{스케일링은 답을 바꾸는 게 아니라 경사하강법의 걸음 수를
    바꾸는 것}.
    \\ 정규방정식은 스케일링 여부와 무관하게 같은 답(단, 스케일이
    바뀌면 계수 스케일도 같이 바뀜)를 주므로, ``답이 바뀌어서''
    스케일링을 하는 게 \emph{아니다}.
  \end{block}
\end{frame}
